% ============================================================================
% skill-translator 中文版期刊论文模板（preamble 部分）
%   用法：build.py 会自动把本文件 + 元数据 + 正文拼成 paper/paper.tex，
%         因此不要在这里写 \begin{document}。
%   字体：documentclass 用 fontset=none，具体的 CJK/西文字体由 build.py 探测生成的
%         paper/fontset.tex 决定（Windows 宋体/黑体、macOS Songti、Linux Noto、
%         TeX Live 自带 Fandol 兜底），以此做到"换一台电脑也能编"。
%   编号：公式沿用原文编号（\tag{n}），图表按译文顺序自动编号；正文引用一律写文字
%         「式(3)」「图1」「表2」，不使用 \eqref 跨语言引用。
% ============================================================================
\documentclass[11pt,a4paper,fontset=none]{ctexart}

\usepackage{fontspec}
\input{fontset.tex}

\usepackage{amsmath,amssymb,amsfonts,bm}
\usepackage{graphicx}
\usepackage{booktabs,array,multirow,tabularx,longtable}
\usepackage{caption}
\usepackage{geometry}
\usepackage{enumitem}
\usepackage{fancyhdr}
\usepackage{xcolor}
\usepackage{indentfirst}
\usepackage{float}
\usepackage{textcomp}
\usepackage[hidelinks]{hyperref}

\geometry{a4paper,left=2.5cm,right=2.5cm,top=2.5cm,bottom=2.5cm}

% ---------- 中文排版常规 ----------
\setlength{\parindent}{2em}
\linespread{1.35}
\setlength{\parskip}{0.15em}

% ---------- 图表题注：中文黑体小号；表题在上、图题在下 ----------
\captionsetup{font={sf,small},labelsep=quad,skip=5pt}
\captionsetup[table]{position=top}
\renewcommand{\figurename}{图}
\renewcommand{\tablename}{表}
\renewcommand{\arraystretch}{1.15}
\setlength{\abovecaptionskip}{5pt}
\setlength{\belowcaptionskip}{4pt}

% ---------- 标题层级：黑体；最多三级 ----------
\ctexset{
  section/format     = \large\bfseries\heiti,
  subsection/format  = \normalsize\bfseries\heiti,
  subsubsection/format = \normalsize\bfseries\heiti,
  paragraph/format   = \normalsize\bfseries\heiti,
}

% ---------- 素材路径：图片在 extracted/figures，表格在 extracted/tables ----------
\graphicspath{{../extracted/figures/}{./figures/}{./assets/}{../extracted/tables/}}
\makeatletter
\def\input@path{{../extracted/tables/}{../extracted/equations/draft/}{./}}
\makeatother

% ---------- 页眉页脚 ----------
\newcommand{\pTwoZHead}{}
\pagestyle{fancy}
\fancyhf{}
\fancyhead[C]{\small\pTwoZHead}
\fancyfoot[C]{\small\thepage}
\renewcommand{\headrulewidth}{0.4pt}
\renewcommand{\footrulewidth}{0pt}

% 公式可跨页；略宽松的断行容忍度（中文+长公式+DOI 链接），避免大量 Overfull
\allowdisplaybreaks
\sloppy

% 论文里常见的写法别名
\providecommand{\eqnref}[1]{式(#1)}
\providecommand{\figref}[1]{图#1}
\providecommand{\tabref}[1]{表#1}
